% Generated draft: requires mathematical and source review.
\documentclass[11pt]{article}
\usepackage[margin=25mm]{geometry}
\usepackage{fontspec}
\setmainfont{Latin Modern Roman}
\usepackage{amsmath,amssymb,mathtools,unicode-math}
\setmathfont{Latin Modern Math}
\usepackage{enumitem,needspace,xurl,hyperref}
\setlength{\parindent}{0pt}
\setlength{\parskip}{5pt}
\setlength{\emergencystretch}{3em}
\newcommand{\N}{\mathbb{N}}
\newcommand{\Z}{\mathbb{Z}}
\newcommand{\Q}{\mathbb{Q}}
\newcommand{\R}{\mathbb{R}}
\newcommand{\C}{\mathbb{C}}
\newcommand{\F}{\mathbb{F}}
\newcommand{\A}{\mathbb{A}}
\newcommand{\PP}{\mathbb{P}}
\newcommand{\E}{\mathbb{E}}
\newcommand{\eps}{\varepsilon}
\newcommand{\dd}{\,\mathrm{d}}
\begin{document}
\section*{UnsolvedMath ID: 701537}
\textbf{Source problem code:} catalogue entry
\subsection{Definitions and mathematical statement}
For an abstract group \(G\), regarded as a discrete group, Kazhdan's property \((T)\) means the following: whenever a unitary representation \(\pi:G\to\mathcal U(H)\) on a complex Hilbert space has almost invariant unit vectors, it has a nonzero invariant vector. Here almost invariant unit vectors means that for every finite subset \(F\subseteq G\) and every \(\varepsilon>0\), there exists \(v\in H\) with \(\|v\|=1\) such that
\[
\|\pi(g)v-v\|<\varepsilon\qquad(g\in F).
\]
An invariant vector satisfies \(\pi(g)v=v\) for every \(g\in G\). A homomorphism onto a group means a surjective group homomorphism.

Let \(\mathcal L_1\) denote the intended class of simple groups of Lie type of rank \(1\), and let \([S]\) denote the isomorphism class of a group \(S\). The proposed assertion is
\[
\text{for every group }G,\qquad
G\text{ has property }(T)
\ \stackrel{?}{\Longrightarrow}\ 
\left|\left\{[S]:S\in\mathcal L_1\text{ and there exists a surjective homomorphism }G\twoheadrightarrow S\right\}\right|<\infty.
\]
Thus the question asks whether a single group with property \((T)\) can have infinitely many pairwise nonisomorphic simple quotients of the indicated type.

\emph{Unresolved qualifications.} The supplied wording does not explicitly restrict the target groups to finite groups or specify the convention for rank. In particular, absolute Lie rank and relative or twisted rank can yield different classes. Consequently, \(\mathcal L_1\) cannot be given an unambiguous membership definition from this record alone. The displayed formulation interprets infinitely many target groups up to isomorphism and uses the discrete-group meaning of property \((T)\); these conventions require confirmation against the original problem.

The record also supplies the following background assertion as known, without a separate proof reference. For a prime power \(q\), let \(\F_q\) be the field with \(q\) elements and put
\[
\operatorname{PSL}_2(q)=\operatorname{SL}_2(\F_q)/Z\bigl(\operatorname{SL}_2(\F_q)\bigr).
\]
The supplied assertion is that, for every group \(G\),
\[
\left|\left\{q:q\text{ is a prime power and }G\twoheadrightarrow\operatorname{PSL}_2(q)\right\}\right|=\infty
\quad\Longrightarrow\quad
G\text{ does not have property }(T).
\]
This background assertion is recorded as source-supplied information, not independently verified here.
\subsection{Short English statement}
Problem 18.109, proposed by A. Jaikin-Zapirain, asks whether every group with Kazhdan's property \((T)\) has only finitely many simple quotients of Lie type of rank \(1\), counted up to isomorphism. The source says the analogous obstruction for quotients \(\operatorname{PSL}_2(q)\) with infinitely many \(q\) is known. The target class and topology conventions need clarification. The registry labels the problem open; this is a supplied status, not an independently verified assessment.
\subsection{Sources}
\begin{itemize}
\item Local UnsolvedMath registry, supplied complete source record: permanent ID \(701537\), exact source problem code \texttt{catalogue entry}, title \emph{Kourovka Notebook Problem 18.109}. The supplied record labels the problem \texttt{open} and records its update as \texttt{2026-10-16T14:00:00Z}.
\item \emph{The Kourovka Notebook (Unsolved Problems in Group Theory)}, Problem 18.109, 18th issue (2014); the supplied background also identifies the 21st issue edition (2026), arXiv identifier \texttt{1401.0300v48}. The supplied collection attribution names E. I. Khukhro and V. D. Mazurov as editors.
\item Proposed by A. Jaikin-Zapirain (2014), according to the supplied record. The assertion concerning infinitely many quotients \(\operatorname{PSL}_2(q)\) is supplied as known, but no separate theorem citation is provided. These bibliographic details and mathematical status statements have not been independently verified.
\end{itemize}
\end{document}
